%==================================================================
% Ini adalah abstrak dalam bahasa indonesia 
%==================================================================

%% DILARANG EDIT BAGIAN INI
\clearpage
\phantomsection
\addcontentsline{toc}{chapter}{ABSTRAK}
\noindent\textbf{\penulis, \nim}\\
\textbf{\MakeUppercase{\judulid}; dibimbing oleh {\pembimbingutama} dan {\pembimbingpendamping}.}\\
% \textbf{Hal.}
\textbf{\pageref{LastPage} + \getromanpagelast\ hal / \total{totaltable} tabel / \total{totalfigure} gambar / \total{lampiran} lampiran / \total{pustaka} pustaka (\tahunakademik)}
\begin{center}
    \textbf{ABSTRAK}\\[0.5cm]
\end{center}
%% DILARANG EDIT BAGIAN INI

%% edit bagian ini
Penyakit pada daun tomat dapat menurunkan kualitas dan produktivitas tanaman apabila tidak dikenali sejak dini. Salah satu pendekatan yang dapat digunakan untuk membantu proses identifikasi penyakit tersebut adalah klasifikasi citra berbasis \textit{deep learning}. Namun, model yang digunakan pada perangkat \textit{edge} tidak hanya perlu memiliki akurasi yang baik, tetapi juga harus ringan dan efisien saat dijalankan pada perangkat dengan sumber daya terbatas. Penelitian ini bertujuan untuk membandingkan performa YOLOv8 Nano dan YOLO11 Nano dalam klasifikasi penyakit daun tomat berbasis perangkat \textit{edge}. Dataset yang digunakan terdiri dari 7.200 citra daun tomat yang terbagi ke dalam enam kelas, yaitu Tomato Early Blight, Tomato Healthy, Tomato Leaf Late Blight, Tomato Leaf Yellow Curl Virus, Tomato Mold Leaf, dan Tomato Septoria Leaf Spot. Model dilatih menggunakan konfigurasi yang sama, kemudian diuji menggunakan 720 citra data uji yang seimbang. Hasil pengujian menunjukkan bahwa YOLOv8 Nano memperoleh akurasi sebesar 88,47\%, \textit{macro F1-score} sebesar 88,63\%, rata-rata waktu inferensi 7,78 ms, dan FPS sebesar 128,52 pada Raspberry Pi 5. Sementara itu, YOLO11 Nano memperoleh akurasi sebesar 88,06\%, \textit{macro F1-score} sebesar 88,12\%, rata-rata waktu inferensi 7,86 ms, dan FPS sebesar 127,18. Berdasarkan hasil tersebut, YOLOv8 Nano menunjukkan performa yang sedikit lebih baik dibandingkan YOLO11 Nano pada skenario penelitian ini. Meskipun selisih performanya kecil, hasil penelitian ini menunjukkan bahwa pengujian langsung pada dataset dan perangkat target tetap diperlukan sebelum menentukan model yang paling sesuai untuk implementasi pada perangkat \textit{edge}.
%% edit sampai sini

%% DILARANG EDIT BAGIAN INI
\noindent\textbf{Kata kunci:} klasifikasi citra, penyakit daun tomat, YOLOv8 Nano, YOLO11 Nano, perangkat \textit{edge}
%% DILARANG EDIT BAGIAN INI
